% GENERATED FILE. Edit canonical lessons, then run npm run generate:books.
% canonical-source-hash: fnv1a64:37f58b44efe4c269
% canonical-lessons: AR-C377-hukuma, AR-C377-wazir, AR-C377-wizara, AR-C377-barlaman, AR-C377-dustur

\chapter{Government}
\label{ch:ar-government}
\hlchaptermodality{377}

\begin{chapteropening}
\textbf{By the end of this chapter:} \emph{I can talk about a government, a minister, a ministry, a parliament and a constitution.}

Complete the last lesson of chapter 377: I can talk about a government, a minister, a ministry, a parliament and a constitution.
\end{chapteropening}

\section[\texorpdfstring{\ar{حكومة} (ḥukūma)}{حكومة (ḥukūma)}]{\ar{حكومة} (ḥukūma) — a government}

\label{lesson:AR-C377-hukuma}



\begin{quote}
Before the new one: say the Arabic for an expert, then the Arabic for academic.
\end{quote}

\subsection*{You'll want to know: \ar{حكومة}}
\textbf{\ar{حكومة}} — \emph{ḥukūma} — \textquotedblleft{}a government\textquotedblright{}.

\begin{quote}
\ar{قررت} \ar{الحكومة} \ar{فتح} \ar{مستشفى} \ar{جديد}.
\end{quote}

\begin{quote}
\emph{qarrarat al-ḥukūma fatḥ mustashfā jadīd.}
\end{quote}

\begin{quote}
The government decided to open a new hospital.
\end{quote}

\begin{cousinweb}
\textbf{\ar{حكومة}} (\emph{ḥukūma}) is built on the native Arabic root \emph{ḥ-k-m}, \textquotedblleft{}to judge, to rule\textquotedblright{} (Lane). \emph{Ḥukūma} is a verbal noun, ruling, which came to mean the body that rules: a government.
\end{cousinweb}

\begin{grammarlens}[title={What you've built}]
One more word for giving reasons and opinions.
\end{grammarlens}

\subsection*{Your turn}
\begin{itemize}
\raggedright
  \item \emph{Say it:} \emph{ḥukūma}
  \item \emph{Say it:} \emph{ḥukūma}, once more
  \item \emph{Say it:} \emph{qarrarat al-ḥukūma fatḥ mustashfā jadīd.} — The government decided to open a new hospital.
  \item \emph{Recall:} say the Arabic for an expert, then the Arabic for academic, then say \emph{ḥukūma} again
\end{itemize}

\begin{tcolorbox}[breakable,colback=teal!4,colframe=teal!35!black,title={Before you move on}]
What does \ar{حكومة} mean? (\textquotedblleft{}A government\textquotedblright{}.) Say it once more.
\end{tcolorbox}
\section[\texorpdfstring{\ar{وزير} (wazīr)}{وزير (wazīr)}]{\ar{وزير} (wazīr) — a minister}

\label{lesson:AR-C377-wazir}



\begin{quote}
Before the new one: say the Arabic for academic, then the Arabic for a government.
\end{quote}

\subsection*{You'll want to know: \ar{وزير}}
\textbf{\ar{وزير}} — \emph{wazīr} — \textquotedblleft{}a minister\textquotedblright{}.

\begin{quote}
\ar{زار} \ar{الوزير} \ar{المدرسة} \ar{أمس}.
\end{quote}

\begin{quote}
\emph{zāra al-wazīr al-madrasa ams.}
\end{quote}

\begin{quote}
The minister visited the school yesterday.
\end{quote}

Plural: \ar{وزراء}, \emph{wuzarāʾ}.

\begin{cousinweb}
\textbf{\ar{وزير}} (\emph{wazīr}) is built on the native Arabic root \emph{w-z-r}, \textquotedblleft{}to bear a heavy load\textquotedblright{} (Lane): the one who carries the ruler's burden. EtymDB lists a Middle Persian cognate, and some trace the word to Persian, so its origin is disputed.
\end{cousinweb}

\begin{grammarlens}[title={What you've built}]
One more word for giving reasons and opinions.
\end{grammarlens}

\subsection*{Your turn}
\begin{itemize}
\raggedright
  \item \emph{Say it:} \emph{wazīr}
  \item \emph{Say it:} \emph{wazīr}, once more
  \item \emph{Say it:} \emph{zāra al-wazīr al-madrasa ams.} — The minister visited the school yesterday.
  \item \emph{Recall:} say the Arabic for academic, then the Arabic for a government, then say \emph{wazīr} again
\end{itemize}

\begin{tcolorbox}[breakable,colback=teal!4,colframe=teal!35!black,title={Before you move on}]
What does \ar{وزير} mean? (\textquotedblleft{}A minister\textquotedblright{}.) Say it once more.
\end{tcolorbox}
\section[\texorpdfstring{\ar{وزارة} (wizāra)}{وزارة (wizāra)}]{\ar{وزارة} (wizāra) — a ministry}

\label{lesson:AR-C377-wizara}



\begin{quote}
Before the new one: say the Arabic for a government, then the Arabic for a minister.
\end{quote}

\subsection*{You'll want to know: \ar{وزارة}}
\textbf{\ar{وزارة}} — \emph{wizāra} — \textquotedblleft{}a ministry\textquotedblright{}.

\begin{quote}
\ar{تعمل} \ar{أمي} \ar{في} \ar{وزارة} \ar{الصحة}.
\end{quote}

\begin{quote}
\emph{taʿmalu ummī fī wizārat aṣ-ṣiḥḥa.}
\end{quote}

\begin{quote}
My mother works at the ministry of health.
\end{quote}

\begin{cousinweb}
\textbf{\ar{وزارة}} (\emph{wizāra}) is built on the native Arabic root \emph{w-z-r}, \textquotedblleft{}to bear a heavy load\textquotedblright{} (Lane). \emph{Wizāra} is the office of a \emph{wazīr}: a ministry.
\end{cousinweb}

\begin{grammarlens}[title={What you've built}]
One more word for giving reasons and opinions.
\end{grammarlens}

\subsection*{Your turn}
\begin{itemize}
\raggedright
  \item \emph{Say it:} \emph{wizāra}
  \item \emph{Say it:} \emph{wizāra}, once more
  \item \emph{Say it:} \emph{taʿmalu ummī fī wizārat aṣ-ṣiḥḥa.} — My mother works at the ministry of health.
  \item \emph{Recall:} say the Arabic for a government, then the Arabic for a minister, then say \emph{wizāra} again
\end{itemize}

\begin{tcolorbox}[breakable,colback=teal!4,colframe=teal!35!black,title={Before you move on}]
What does \ar{وزارة} mean? (\textquotedblleft{}A ministry\textquotedblright{}.) Say it once more.
\end{tcolorbox}
\section[\texorpdfstring{\ar{برلمان} (barlamān)}{برلمان (barlamān)}]{\ar{برلمان} (barlamān) — a parliament}

\label{lesson:AR-C377-barlaman}



\begin{quote}
Before the new one: say the Arabic for a minister, then the Arabic for a ministry.
\end{quote}

\subsection*{You'll want to know: \ar{برلمان}}
\textbf{\ar{برلمان}} — \emph{barlamān} — \textquotedblleft{}a parliament\textquotedblright{}.

\begin{quote}
\ar{ناقش} \ar{البرلمان} \ar{خطة} \ar{الحكومة}.
\end{quote}

\begin{quote}
\emph{nāqasha al-barlamān khiṭṭat al-ḥukūma.}
\end{quote}

\begin{quote}
Parliament discussed the government's plan.
\end{quote}

\begin{cousinweb}
\textbf{\ar{برلمان}} (\emph{barlamān}) was borrowed from French \emph{parlement} (EtymDB), a place of speaking. Arabic has no p, so the word begins with b.
\end{cousinweb}

\begin{grammarlens}[title={What you've built}]
One more word for giving reasons and opinions.
\end{grammarlens}

\subsection*{Your turn}
\begin{itemize}
\raggedright
  \item \emph{Say it:} \emph{barlamān}
  \item \emph{Say it:} \emph{barlamān}, once more
  \item \emph{Say it:} \emph{nāqasha al-barlamān khiṭṭat al-ḥukūma.} — Parliament discussed the government's plan.
  \item \emph{Recall:} say the Arabic for a minister, then the Arabic for a ministry, then say \emph{barlamān} again
\end{itemize}

\begin{tcolorbox}[breakable,colback=teal!4,colframe=teal!35!black,title={Before you move on}]
What does \ar{برلمان} mean? (\textquotedblleft{}A parliament\textquotedblright{}.) Say it once more.
\end{tcolorbox}
\section[\texorpdfstring{\ar{دستور} (dustūr)}{دستور (dustūr)}]{\ar{دستور} (dustūr) — a constitution}

\label{lesson:AR-C377-dustur}



\begin{quote}
Before the new one: say the Arabic for a ministry, then the Arabic for a parliament.
\end{quote}

\subsection*{You'll want to know: \ar{دستور}}
\textbf{\ar{دستور}} — \emph{dustūr} — \textquotedblleft{}a constitution\textquotedblright{}.

\begin{quote}
\ar{كتب} \ar{الشعب} \ar{دستورا} \ar{جديدا}.
\end{quote}

\begin{quote}
\emph{kataba ash-shaʿb dustūran jadīdan.}
\end{quote}

\begin{quote}
The people wrote a new constitution.
\end{quote}

Plural: \ar{دساتير}, \emph{dasātīr}.

\begin{cousinweb}
\textbf{\ar{دستور}} (\emph{dustūr}) was borrowed from Persian \emph{dastūr} (EtymDB); Lane gives it as an arabicized word for the register of a realm's rules (Lane). Modern Arabic made it the constitution.
\end{cousinweb}

\begin{grammarlens}[title={What you've built}]
One more word for giving reasons and opinions.
\end{grammarlens}

\subsection*{Your turn}
\begin{itemize}
\raggedright
  \item \emph{Say it:} \emph{dustūr}
  \item \emph{Say it:} \emph{dustūr}, once more
  \item \emph{Say it:} \emph{kataba ash-shaʿb dustūran jadīdan.} — The people wrote a new constitution.
  \item \emph{Recall:} say the Arabic for a ministry, then the Arabic for a parliament, then say \emph{dustūr} again
\end{itemize}

\begin{tcolorbox}[breakable,colback=teal!4,colframe=teal!35!black,title={Before you move on}]
What does \ar{دستور} mean? (\textquotedblleft{}A constitution\textquotedblright{}.) Say it once more.
\end{tcolorbox}
